% \iffalse meta-comment
%
% hit-thesis-titlepage.dtx 是内封（扉页）
%
% \fi
%
% \section{内封}
%
%    \begin{macrocode}
%<*thesis-titlepage>
%    \end{macrocode}
%
% \changes{v3.1d}{2025/03/03}{增加哈尔滨本科的第二封面。\issue[172]}
% \changes{v3.2a}{2026/08/16}{本科内封三校区统一，威海那一套删掉，
%   \cs{hit@second@title@page@harbin@bachelor} 改名 \cs{hit@second@title@page@bachelor}}
% 继深圳之后，本部的本科模板又开始整新活了，大改一通，醉好的结局就是都统一吧。
% 现在统一了。
%
% 第二张封面有三套版式：博后、本科、其余。
%    \begin{macrocode}
\cs_new_protected:Npn \hit@titlepage
  {
    \bool_if:NTF \g__hit_postdoc_bool
      { \hit@titlepage@postdoc }
      {
        \bool_if:NTF \g__hit_bachelor_bool
          { \hit@titlepage@bachelor }
          { \hit@titlepage@other }
      }
  }
%    \end{macrocode}
%
%    \begin{macrocode}
\cs_new_protected:Npn \hit@titlepage@postdoc
  {
~
 \xiaosi[1]\songti\textcolor[rgb]{1,1,1}{{\hit@hi}}%
 \vspace*{1.2cm}~
 \begin{center}~
 \parbox[t][8cm][t]{\textwidth}{~
 \begin{center}\songti\sanhao[2]\hit@title@cover@zh\end{center}~
 \vspace*{16bp}~
 \begin{center}\sihao[2]\MakeUppercase{\hit@title@cover@en}\end{center}~
 }~
%    \end{macrocode}
%
% \changes{v3.2a}{2026/08/08}{博后封面那两张字段表改用 \cs{makebox} 定宽，
%   不再借 \pkg{tabularx}。两张表要左边缘对齐，靠的是同一个 0.6\cs{linewidth}；
%   \pkg{tabularx} 没有 X 列时也会把盒子撑到指定宽度，效果一样但每张表报一条
%   Overfull，日志里两条噪音}
% 两张表的第一列宽度不同（上表标签撑到 7 个字，下表是八个字的「研究工作起始时间」），
% 各自按自然宽度居中就会错开。所以两张都装进同宽的 \cs{makebox} 里左对齐，
% 左边缘落在同一处。
%    \begin{macrocode}
 \parbox[t][5cm][t]{\textwidth}{~
 \begin{center}\songti\sihao[2]~
 \makebox[0.6\linewidth][l]{\begin{tabular}{l@{\hit@postdoc@separator}l}~
 \hit@postdoc@author@field@label &~\hit@author@zh \\~
 \hit@postdoc@co@supervisor@field@label &~\hit@supervisor@zh \\~
 \hit@postdoc@station@field@label &~\hit@station@zh \\~
 \hit@postdoc@major@field@label &~\hit@major@zh \\~
 \end{tabular}}~
 \makebox[0.6\linewidth][l]{\begin{tabular}{l@{\hit@postdoc@separator}l}~
 \hit@postdoc@research@start@date@field@label &~\hit@research@start@date@zh \\~
 \hit@postdoc@research@end@date@field@label &~\hit@research@end@date@zh \\~
 \end{tabular}}
 \end{center}
 }
 \end{center}~
 \vfill
 {\songti\sihao[2]~
 \begin{center}~
 \hit@school@zh%
 （\hit@campus）%
 \end{center}~
 \begin{center}~
 \hit@date@zh
 \end{center}}
  }
%    \end{macrocode}
% \begin{macro}{\hit@titlepage@line}
% \changes{v3.2a}{2026/08/08}{第二封面上“字段名加值”的行型抽成
%   \cs{hit@titlepage@line}}
% 第二封面的信息表不画下划线，字段名走黑体。哈尔滨本科那张七行、其余那张九行，
% 行型完全一样。带星号的不出 \cs{\textbackslash\textbackslash}，给最后一行用。
%    \begin{macrocode}
\NewDocumentCommand \hit@titlepage@line { s m m }
  { {\heiti \hit@put@title{#2}}~&~#3 \IfBooleanF {#1} { \\ ~ } }
%    \end{macrocode}
% \end{macro}
%    \begin{macro}{\hit@empty@box}
%    \begin{macrocode}
\cs_new_protected:Npn \hit@empty@box
  {
    {
      \dim_set:Nn \fboxsep { 0pt } \dim_set:Nn \fboxrule { 0.5pt }
      \raisebox { 0.5pt } { \framebox [8pt] [c] { \raisebox { 0pt } [7pt] [0pt] { } } }
    }
  }
%    \end{macrocode}
%    \end{macro}
%    \begin{macro}{\hit@checked@box}
%    \begin{macrocode}
\cs_new_protected:Npn \hit@checked@box
  {
    {
      \dim_set:Nn \fboxsep { 0pt } \dim_set:Nn \fboxrule { 0.5pt }
      \raisebox { 0.5pt }
        {
          \framebox [8pt] [c]
            {
              \raisebox { 0pt } [7pt] [0pt]
                {
                  \makebox [0pt] [c]
                    { \raisebox { 1.2pt } { \resizebox { 6.5pt } { 6.5pt } { $ \checkmark $ } } }
                }
            }
        }
    }
  }
%    \end{macrocode}
%    \end{macro}
%    \begin{macro}{\hit@titlepage@bachelor}
% \changes{v3.2a}{2026/08/07}{本科生内封上方文字添加毕业论文或设计的选择框，由于目前hitheis不支持毕业设计，因此对勾写死在毕业论文上}
% \changes{v3.2a}{2026/08/07}{schrb: 同步为哈工深官方设计的样式}
%    \begin{macrocode}
%    \end{macrocode}
%
% \changes{v3.2a}{2026/08/24}{本科内封改照 Word 范例的段落序列排，间隔段与字段
%   行距都从 docx 的 \texttt{w:pPr} 读出来}
%
% 内封那一节与封面同档，\texttt{docGrid} 的 \texttt{linePitch} 是 395 缇。
% 序列取值（本科书写范例第三页，批注段换成同高空段）：
% \begin{itemize}
%   \item 标签，宋体小二加粗，不贴网格
%   \item 中文题目，黑体二号，\emph{贴网格}
%   \item 题目与字段之间三个间隔段，\texttt{w:line="780" w:lineRule="exact"}，
%         即固定值 39\,bp
%   \item 字段表，定宽三列 1806 / 301 / 2951 缇，居中、无单元格边距，
%         行高最小 540 缇；每行 \texttt{w:line="360"} 即 1.5 倍，且贴网格，
%         所以行距是 $1.5 \times 19.75 = 29.625$，不是 $1.5 \times$ 自然行高
% \end{itemize}
%
% 范例那一页第一行是「密级：公开」，我们这一版第一行还带毕业论文与毕业设计的
% 勾选框，两份模板不同版，所以绝对基线不对拍，对的是相对间距。
%    \begin{macrocode}
\tl_const:Nn \c__hit_titlepage_layout_tl { line-pitch = 19.75bp , }
\dim_const:Nn \c__hit_titlepage_label_dim { 90.3bp }
\dim_const:Nn \c__hit_titlepage_colon_dim { 15.05bp }
\dim_const:Nn \c__hit_titlepage_value_dim { 147.55bp }
\cs_new_protected:Npn \hit@titlepage@field:nn #1#2
  {
    \hit@format@par:nn { titlepage-field }
      {
        \hbox_to_wd:nn { \c__hit_titlepage_label_dim }
          { \hit@put@title {#1} \hfil }
        \hbox_to_wd:nn { \c__hit_titlepage_colon_dim }
          { \hit@title@separator@zh \hfil }
        \hbox_to_wd:nn { \c__hit_titlepage_value_dim } { #2 \hfil }
      }
  }
\cs_new_protected:Npn \hit@titlepage@bachelor
  {
    \exp_args:Nx \hit@msword@layout@push:n
      {
        \exp_not:V \g__hit_layout_msword_tl
        \exp_not:V \c__hit_titlepage_layout_tl
      }
    \group_begin:
      \dim_zero:N \parindent
      \dim_zero:N \parskip
      \hit@format@par:nn { titlepage-mark }
        {
          \hbox_to_wd:nn { \textwidth }
            {
              \hit@if@option@TF { design }
                { \hit@empty@box } { \hit@checked@box }
              \nobreakspace 毕业论文 \quad
              \hit@if@option@TF { design }
                { \hit@checked@box } { \hit@empty@box }
              \nobreakspace 毕业设计
              \hfil
              \hit@secret@level@field@label \hit@title@separator@zh
              \hit@secret@level
            }
        }
      \hitenter*[lines=4,size=xiaosi,font-zh=songti,line-spacing=1]
      \hit@format@par:nn { titlepage-label }
        { \hit@bachelor@degree@level \hit@bachelor@document@type }
      \hitenter*[lines=1,size=xiaosi,font-zh=songti,line-spacing=19.75bp]
      \hitenter*[lines=1,size=xiaosi,font-zh=songti,line-spacing=1.25]
      \hit@format@par:nn { titlepage-title } { \hit@title@cover@zh }
      \hit@if@subtitle@T
        {
          \hit@format@par:nnn { titlepage-title } { size = xiaoer }
            { ——\hit@title@subtitle@zh }
        }
      \hitenter*[lines=2,size=xiaosi,font-zh=songti,line-spacing=19.75bp]
      \hitenter*[lines=3,size=erhao,font-zh=songti,line-spacing=39bp]
      \dim_set_eq:NN \hit@title@width \c__hit_titlepage_label_dim
      \hit@titlepage@field:nn
        { \hit@author@field@label@zh } { \hit@author@zh }
      \hit@titlepage@field:nn
        { \hit@bachelor@student@id@field@label } { \hit@student@id }
      \hit@titlepage@field:nn
        { \hit@bachelor@supervisor@field@label } { \hit@supervisor@zh }
      \hit@titlepage@field:nn
        { \hit@bachelor@major@field@label } { \hit@major@zh }
      \hit@titlepage@field:nn
        { \hit@bachelor@affiliation@field@label } { \hit@affiliation@zh }
      \hit@titlepage@field:nn
        { \hit@date@field@label@zh } { \hit@date@zh }
      \hit@titlepage@field:nn
        { \hit@school@field@label@zh } { \hit@school@zh }
    \group_end:
    \hit@msword@layout@pop:
  }
%    \end{macrocode}
%    \end{macro}
%    \begin{macro}{\hit@titlepage@other}
%    \begin{macrocode}

%内封
\cs_new_protected:Npn \hit@titlepage@other
  {
~
 \begin{center}~
%    \end{macrocode}
%
% 生成头部信息，使用 \cs{hfill} 将第二列信息推到右侧。
% \changes{v3.1b}{2022/06/06}{修改本部硕士第二封面的信息的对齐方式}
% 本部硕士第二封面的信息的对齐方式不同。
%    \begin{macrocode}
 {~
 \songti \xiaosi
 \hit@if@both@options@TF{harbin}{master}~
 {~
 \begin{tabular}{l}~
 \hit@classified@index@national@field@label@zh：~\hit@classified@index@national\\~
 \hit@classified@index@international@field@label@zh：~\hit@classified@index@international
 \end{tabular}\hfill
 \begin{tabular}{l}~
 \hit@school@id@field@label：\hit@school@id\\~
 \hit@secret@level@field@label：\hit@secret@level
 \end{tabular}
 }%
 {~
 \begin{tabular}{@{}r@{：}l@{}}~
 \hit@classified@index@national@field@label@zh &~\hit@classified@index@national\\~
 \hit@classified@index@international@field@label@zh &~\hit@classified@index@international
 \end{tabular}\hfill
 \begin{tabular}{@{}r@{：}l@{}}~
 \hit@school@id@field@label &~\hit@school@id\\~
 \hit@secret@level@field@label &~\hit@secret@level
 \end{tabular}
 }%
 }~
 \parbox[t][3.2cm][t]{\textwidth}{\begin{center}~\end{center}~}~
%    \end{macrocode}
%
% \changes{v3.1b}{2022/06/06}{本部硕/博士第二页““”学硕/博士学位论文”中删除““”学”}
% \changes{v3.1c}{2023/04/16}{修正本部硕(博)士(这里的空格)学位论文}
% \changes{v3.1b}{2022/06/06}{深圳硕士第二页““”学硕士学位论文”中删除““”学”}
% 本部硕/博士第二页设置为“硕/博士学位论文”。深圳的我把博士也一块改了，要不然要重写一下判断，如果发现博士不需要改的话再重写吧…
% 博士在 \issue[222] 还是需要改...
% \changes{v3.1d}{2025/03/03}{修改深圳博士第二封面，\issue[222]}
%    \begin{macrocode}
 \parbox[t][2.4cm][t]{\textwidth}{~
 \xiaoer[1]~
 \begin{center}~
 \songti\bfseries
 \hit@if@harbin@master@or@doctor@or@shenzhen@master@TF{\hit@degree@level@zh}{\hit@degree@of@discipline@format@zh}\hit@document@type
 \end{center}}~
 \parbox[t][5cm][t]{\textwidth}{\erhao
 \begin{center}\heiti\hit@title@cover@zh\end{center}~
\hit@if@subtitle@T{\begin{center}\hspace{-4em}\xiaoer\heiti——\hit@title@subtitle@zh\end{center}}}~
 \parbox[t][9.8cm][b]{\textwidth}~
 {\sihao
 \dim_set:Nn \hit@title@width {6em}~
 \begin{center}~\cs_set:Npn \arraystretch {1.62}~\songti
 \begin{tabular}{l@{\hit@title@separator@zh}l}~
 \hit@titlepage@line{\hit@author@field@label@zh}{\hit@author@zh}
 \hit@titlepage@line{\hit@supervisor@field@label@zh}{\hit@supervisor@zh}
 \tl_if_empty:NF\hit@associate@supervisor@zh{%
 {\heiti \hit@put@title{\hit@associate@supervisor@field@label@zh}}&~\hit@associate@supervisor@zh\\~
 }%
 \tl_if_empty:NF\hit@co@supervisor@zh{%
 \hit@titlepage@line{\hit@co@supervisor@field@label@zh}{\hit@co@supervisor@zh}
 }%
 \hit@titlepage@line{\hit@degree@field@label@zh}{\hit@degree@of@discipline@format@zh}
 \hit@titlepage@line{\hit@major@field@label@zh}{\hit@major@zh}
 \hit@titlepage@line{\hit@affiliation@field@label@zh}{\hit@affiliation@zh}
 \hit@titlepage@line{\hit@date@field@label@zh}{\hit@date@zh}
 \hit@titlepage@line*{\hit@school@field@label@zh}{\hit@school@zh}
 \end{tabular}\cs_set:Npn \arraystretch {1}
 \end{center}~}
 \end{center}~
  }

%    \end{macrocode}
%    \end{macro}
%    \begin{macro}{\hit@titlepage@en}
%    \begin{macrocode}

% 英文封面
\NewDocumentCommand \emultiline { O{c} m }
  {
    \cs_set:Npn \arraystretch { 1 }
    \begin { tabular } [#1] { @{}l@{} } #2 \end { tabular }~
    \cs_set:Npn \arraystretch { 1.3 }
  }
\cs_new_protected:Npn \hit@titlepage@en
  {
~
 {%
 \xiaosi[1.667]\noindent {\hit@classified@index@national@field@label@en}:~\hit@classified@index@national \\[8pt]~
 {\hit@classified@index@international@field@label@en}:~\hit@classified@index@international %
 }~
 \vspace*{1em}~
 \begin{center}~
 %    \end{macrocode}
%
% 深圳硕士英文封面第一行文字不同，修改。又要修改啦，哈哈。
% 本部硕博英文封面第一行不包含学科类别。\issue[123]
% \changes{v3.0.09}{2020/10/17}{深圳硕士英文封面第一行文字不同，修改}
% \changes{v3.1b}{2022/06/06}{修改本部硕博英文封面第一行。\issue[123]}
% \changes{v3.1c}{2023/04/16}{深圳硕士模板会同时出现两个英文标题}
% \changes{v3.1c}{2023/04/16}{修改深圳硕士模板的 Engineering 为 \cs{hit@exueke}}
% \changes{v3.1c}{2023/04/16}{按照 \issue[97] 修改深圳硕士的英文封面}
%    \begin{macrocode}
 \parbox[t][1.6cm][t]{\textwidth}{\begin{center}~\end{center}~}~
 \parbox[t][3.5cm][t]{\textwidth}{\xiaoer[1]~
 \begin{center}~
 \hit@if@both@options@TF{shenzhen}{master}~
 {%
 Dissertation~for~the~Master~Degree~
 }~
 {%
 \hit@if@harbin@master@or@doctor@TF
 {Dissertation~for~the~{\hit@degree@level@attribute}~Degree}~
 {Dissertation~for~the~{\hit@degree@level@attribute}~Degree~in~\hit@discipline@en}
 }
 \end{center}~
 \hit@if@option@TF{fulltime}{}%
 {\begin{center}\hit@brace@left@en\hit@student@type@text:n{en}\hit@brace@right@en\end{center}}%
 }~%与中文保持一致，删除in {\hit@discipline@en}
 \parbox[t][7cm][t]{\textwidth}{%
%    \end{macrocode}
%
% \changes{v3.1d}{2025/03/03}{根据 \issue[227] 修改英文标题不能换行的问题}
%    \begin{macrocode}
\skip_if_exist:NF \hit@etitlelength { \skip_new:N \hit@etitlelength }
 \settowidth{\hit@etitlelength}{\erhao\hit@title@cover@en\hit@if@subtitle@T{\hit@title@separator@en\hit@title@subtitle@en}}%
 \begin{center}%
% \changes{v3.2a}{2026/08/07}{英文内封标题默认二号字，可通过etitleauto选项启用自动调整}
 \hit@if@option@TF{etitleauto}%
 {\dim_compare:nNnTF{\hit@etitlelength}>{450mm}{\xiaoer}{\erhao}}%
 {\erhao}%
 \textbf{\MakeUppercase{\hit@title@cover@en}%
\hit@if@subtitle@T{\hit@title@separator@en\MakeUppercase{\hit@title@subtitle@en}}}\end{center}}~
 %★★★★若信息内容不太长，不会引起信息内容分行时，使用tabular环境，否则使用下面的tabularx环境。
 {\sihao\cs_set:Npn \arraystretch {1.3}~
 \begin{tabular}{@{}l@{\nobreakspace{}}l@{}}~
 \textbf{\hit@author@field@label@en\hit@title@separator@en}~&~\hit@author@en\\~
 \textbf{\hit@supervisor@field@label@en\hit@title@separator@en}~&~\hit@supervisor@en\\~
 \tl_if_empty:NF\hit@associate@supervisor@en{%
 \textbf{\hit@associate@supervisor@field@label@en\hit@title@separator@en}~&~\hit@associate@supervisor@en\\~
 }%
 \tl_if_empty:NF\hit@co@supervisor@en{%
 \textbf{\hit@co@supervisor@field@label@en\hit@title@separator@en}~&~\hit@co@supervisor@en\\~
 }%
 \textbf{\hit@degree@field@label@en\hit@title@separator@en}~&~\hit@degree@of@discipline@format@en\\~
 \textbf{\hit@major@field@label@en\hit@title@separator@en}~&\hit@major@en\\~
 \textbf{\hit@affiliation@field@label@en\hit@title@separator@en}~&\hit@affiliation@en\\~
 \textbf{\hit@date@field@label@en\hit@title@separator@en}~&~\hit@date@en\\~
 \textbf{\hit@school@field@label@en\hit@title@separator@en}~&~\hit@school@en
 \end{tabular}\cs_set:Npn \arraystretch {1}}
 \end{center}~
  }

%    \end{macrocode}
%    \end{macro}
%
%    \begin{macrocode}
%</thesis-titlepage>
%    \end{macrocode}
%
% \Finale
